\scanpage{12}
% L01–03
``Main Theorem'' de Zariski) et surjectif ; et il suffit de prouver la $=$ chose pour $\mathbb E^n\xrightarrow{A_*}\mathbb E^n$, donc prouver que sur $k$ corps alg. clos, l'application
% L04
\deleted{\unsure{a)|2)} \uline{Surjectivité}}
\[ k^{\unsure{*|n}}\xrightarrow{A_*}k^n \]
% L05
est \uline{surjective} à \uline{fibres finies}.

% L06–08
1) Surjection. Soit $a=(a_0,\text{--},a_{n-1})\in k^n$. Si $a=0$, c'est l'image de $0=(\underbrace{0,\text{--},0}_{n})\in k^n$. \deleted{\unsure{Si} $\longleftarrow$}
% L08–11
Sinon, soit $a_r$ le premier parmi $a_0\text{--}a_{\unsure{n|n-1}}$ qui soit $\neq0$. Considérons $a_r\text{--}a_n\in k^{n-r+1}$, il provient de $\xi_1,\text{--},\xi_{n-r+1}$ ; on complète par des $0$ pour avoir $\xi_1\text{---}\xi_n$, ça marche, grâce au
% L12–13
\uline{Lemme} \[
\left\{\begin{aligned}
 A_0(X_1,\text{--},X_{n-1},0)&=0\\
 A_i(\deleted{\unsure{0,}}X_1,\text{--},X_{n-1},0)&=A_{i-1}(X_1,\text{--},X_{n-1})
\end{aligned}\right.
\]
pour $1\leq i\leq\unsure{n|n-1}$
% L14–15
donc
\[
 A_\alpha[X_1,\text{--},X_{n-r},\overbrace{0\text{--},0}^{\unsure{r|n-r|s}}]
 =\begin{cases}
 0&\text{si }0\leq\alpha\unsure{\leq r-1|<r|<s}\\
 A_{\alpha-r}(X_1,\text{--},X_{n-r})&\text{si }\unsure{r|s}\leq\alpha\leq n
 \end{cases}
\]

% L16–17
2) À fibres finies. En fait, on voit \unsure{tt|lt} de suite que les fibres sont les trajectoires de $\Gamma_n$.

% L18–20
[Ainsi, indépendamment du M. Th. de Zar, on voit que $Q_n\to\mathbb E^n$ est surjectif, radiciel, birationnel.

% L21–23
\uline{NB} On peut aussi regarder les invariants de $\mathbb F_p[X_1,\text{--},X_n]$ sous $\unsure{S\Gamma_n|\Sigma\Gamma_n}=\operatorname{Sl}(n,\mathbb F_p)$. Parmi ces invariants il y a $\Delta_n$ (on a $\deleted{\unsure{\Delta|\triangle}}$ en fait
% archive sheet number: 11
